\documentclass[11pt,a4paper]{ctexart}
\usepackage[margin=2.2cm]{geometry}
\usepackage{amsmath,amssymb,amsthm,mathtools,booktabs,hyperref,microtype}
\hypersetup{colorlinks=true,linkcolor=black,urlcolor=blue,citecolor=black}
\newtheorem{theorem}{定理}
\newtheorem{lemma}{引理}
\newtheorem{corollary}{推论}
\renewcommand{\proofname}{证明}
\newcommand{\Fthree}{F_3}
\newcommand{\vthree}{v_3}
\title{\bfseries 平移乘积过程的双根坐标、\\联合深度律与可逆重建}
\author{研究札记 · F3-ROOT}
\date{2026 年 9 月 29 日}
\begin{document}
\maketitle
\begin{abstract}
本文研究配对乘积
\[
a_j=n+2j3^k d,\qquad
b_j=n+(2j3^k+2)d,\qquad
D_j=F_3(a_jb_j).
\]
在前两组进入深层的分支上，归一化后的 \(Q_j=a_jb_j+1\) 是模 \(3\) 等于 \(j(j-1)\) 的首一二次多项式，因此具有两个唯一提升的简单三进根。整条无限深度序列被两个根精确控制，形成两条可升层分支与一条恒定分支。两个根在条件 Haar 测度下独立均匀，由此得到所有有限门槛事件的闭式联合律；反过来，完整深度序列可恢复两个根并显式重建原参数。把结构扩展到六个素数位置时出现强制的模 \(5\) 障碍，令 \(d=5r\) 后可恢复七素数全局实现，并得到模式 \((3,4,2)\) 的渐近比例 \(4/729\)。
\end{abstract}

\section{二次平移乘积}
固定 \(k\ge1\)，令
\[
q=3^k.
\]
定义
\[
a_j=n+2jqd,\qquad
b_j=n+(2jq+2)d,
\]
\[
Q_j=a_jb_j+1.
\]
展开得
\[
\boxed{
Q_j=4q^2d^2j^2+4qd(n+d)j+Q_0.
}
\]
所以
\[
Q_{j+2}-2Q_{j+1}+Q_j=8q^2d^2.
\]

定义基础分支
\[
E_k:\quad
d\equiv1\pmod3,\qquad
v_3(Q_0),v_3(Q_1)\ge2k+1.
\]
在 \(E_k\) 上自动有 \(n\equiv2\pmod3\)，从而
\[
a_j\equiv2,\qquad b_j\equiv1\pmod3,
\]
因此
\[
D_j:=F_3(a_jb_j)=v_3(Q_j).
\]

由于 \(v_3(8q^2d^2)=2k\)，而 \(Q_0,Q_1\) 的赋值至少为 \(2k+1\)，二阶差分立即给出
\[
\boxed{D_2=2k.}
\]

\section{双根因子分解}
定义
\[
P(X)=\frac{Q_X}{4q^2d^2}
=X^2+BX+C,
\]
其中
\[
B=\frac{n+d}{qd},
\qquad
C=\frac{Q_0}{4q^2d^2}.
\]
由
\[
Q_1-Q_0=4qd(n+d+qd)
\]
和 \(E_k\) 得
\[
B\equiv-1\pmod3,\qquad C\equiv0\pmod3.
\]
所以
\[
P(X)\equiv X(X-1)\pmod3.
\]
模 \(3\) 的两个根 \(0,1\) 都是简单根，Hensel 提升给出唯一
\[
r_0\in3\mathbb Z_3,\qquad
r_1\in1+3\mathbb Z_3.
\]
于是
\[
\boxed{
Q_j=4q^2d^2(j-r_0)(j-r_1)
}
\]
以及
\[
\boxed{
D_j=2k+v_3(j-r_0)+v_3(j-r_1).
}
\]

写
\[
r_0=3Z_0,\qquad r_1=1+3Z_1,
\]
得到
\[
\boxed{
\begin{aligned}
D_{3u}&=2k+1+v_3(u-Z_0),\\
D_{3u+1}&=2k+1+v_3(u-Z_1),\\
D_{3u+2}&=2k.
\end{aligned}}
\]

\section{根坐标与参数反演}
令
\[
\sigma=r_0+r_1,\qquad
\delta=r_1-r_0.
\]
比较二次多项式线性系数：
\[
\frac{n+d}{qd}=-\sigma,
\]
故
\[
\boxed{n=-d(1+q\sigma).}
\]
判别式满足
\[
\delta^2
=
\frac{(n+d)^2-(n(n+2d)+1)}{q^2d^2}
=
\frac{d^2-1}{q^2d^2}.
\]
因此
\[
\boxed{
d=(1-q^2\delta^2)^{-1/2},
}
\]
取唯一的 \(d\equiv1\pmod3\) 分支。

反过来，任意
\[
r_0\in3\mathbb Z_3,\qquad
r_1\in1+3\mathbb Z_3
\]
都由这两个公式给出唯一的 \(E_k\) 参数。因此
\[
E_k\longleftrightarrow\mathbb Z_3^2
\]
具有显式双射坐标。

\section{独立 Haar 根与联合尾分布}
令
\[
S=Z_0+Z_1,\qquad
\Delta=Z_1-Z_0.
\]
则
\[
\sigma=1+3S,\qquad
\delta=1+3\Delta.
\]
由反演公式
\[
\frac{\partial d}{\partial\Delta}
=
3q^2\delta d^3,
\]
故其赋值为 \(2k+1\)；而
\[
\frac{\partial n}{\partial S}
=
-3qd
\]
的赋值为 \(k+1\)。

坐标变换 \((Z_0,Z_1)\mapsto(S,\Delta)\) 的行列式为单位 \(2\)，因此总 Jacobian 的赋值恒为
\[
3k+2.
\]
所以条件于 \(E_k\) 后，\((Z_0,Z_1)\) 正是 \(\mathbb Z_3^2\) 上的乘积 Haar 坐标；同时
\[
\mu(E_k)=3^{-(3k+2)}.
\]

令 \(L=2k+1\)。则
\[
D_{3u}\ge L+t
\iff
Z_0\equiv u\pmod{3^t},
\]
\[
D_{3u+1}\ge L+t
\iff
Z_1\equiv u\pmod{3^t}.
\]
同一根上的有限球条件要么不相容，要么等于最深的那个球。若两个分支内部均相容，所有门槛联合概率为
\[
\boxed{
3^{-(T_0+T_1)},
}
\]
其中 \(T_c\) 是第 \(c\) 个根分支上的最大门槛深度；若不相容则概率为 \(0\)。

\section{完整区段层级计数}
对 \(1\le t\le T\)，
\[
D_j\ge2k+t
\]
等价于
\[
j\equiv r_0\quad\text{或}\quad r_1\pmod{3^t}.
\]
两个类不相交，所以任意连续 \(3^T\) 个索引中
\[
\boxed{
\#\{j:D_j\ge2k+t\}
=
2\cdot3^{T-t}.
}
\]
特别地，连续九项中恰有三项等于 \(2k\)、四项等于 \(2k+1\)、两项至少为 \(2k+2\)。

\section{深度序列可逆恢复参数}
在
\[
0\le j<3^T
\]
中，满足
\[
D_j\ge2k+T
\]
的恰有两个位置，它们分别给出
\[
r_0\bmod3^T,\qquad r_1\bmod3^T.
\]
模 \(3\) 余数 \(0,1\) 区分两根。

因此整条 \((D_j)_{j\ge0}\) 恢复 \(r_0,r_1\)，再由反演公式恢复 \(n,d\)。

若根已知模 \(3^T\)，则
\[
d\bmod3^{2k+T},
\qquad
n\bmod3^{k+T}
\]
均可恢复。

逐位恢复时，每个根的下一位有三个候选，测试两个后第三个可由排除确定，因此从模 \(3\) 提升到模 \(3^T\) 最多需要
\[
\boxed{4(T-1)}
\]
次自适应阈值查询。

\section{六位置的模 \(5\) 障碍}
取 \(k=1\)。前三组六个位置是
\[
n,\ n+2d,\ n+6d,\ n+8d,\ n+12d,\ n+14d.
\]
系数
\[
\{0,2,6,8,12,14\}
\]
模 \(5\) 覆盖全部五个剩余类。若 \(5\nmid d\)，乘以 \(d\) 是模 \(5\) 置换，所以六个位置中必有一个被 \(5\) 整除。

若六个位置都大于 \(5\) 且均为素数，则必有
\[
\boxed{5\mid d.}
\]
而 \(E_1\) 又要求 \(d\equiv1\pmod3\)，所以 \(d\) 本身不能是素数。

\section{七素数修复与 \(4/729\)}
令
\[
d=5r.
\]
考虑
\[
r,\ n,\ n+10r,\ n+30r,\ n+40r,\ n+60r,\ n+70r.
\]
这些方向两两不同。

模 \(2\)、模 \(5\) 均显然可容许；模 \(3\) 的允许类恰为
\[
(n,r)\equiv(1,1),(2,2),
\]
占 \(2/9\)。对 \(\ell\ge7\)，固定 \(r\ne0\) 后六个端点最多排除六个 \(n/r\) 类，而 \(6<\ell\)。所以奇异乘积严格为正。

由一般有限复杂度素数线性形式定理，
\[
\#\mathcal P(X)
=
(\mathfrak S+o(1))
\frac{X^2}{(\log X)^7},
\qquad \mathfrak S>0.
\]

在 \(k=1\) 时
\[
\mu(E_1)=3^{-5}.
\]
条件于 \(E_1\)，
\[
D_0=3
\]
的概率为 \(2/3\)，
\[
D_1=4
\]
的概率为 \(2/9\)，且两个根独立；\(D_2=2\) 自动成立。所以精确事件在 $d\equiv1\pmod3$ 分支的绝对三进测度为
\[
3^{-5}\cdot\frac23\cdot\frac29
=
\frac4{3^8}.
\]
同时取负 $(n,d)\mapsto(-n,-d)$ 保持每个 $Q_j$，并在有限三进剩余类上交换两个非零 $d$ 分支。因此完整事件还有等量的第二分支，总测度为 $8/3^8$；这里不将正素数映成负素数，只使用参数剩余类的保测度双射。$d\equiv0\pmod3$ 不可能满足事件。

除以整个七素数系统模 \(3\) 的允许域 \(2/9\)，得到
\[
\boxed{
\Pr((D_0,D_1,D_2)=(3,4,2)\mid\text{七形式全素})
\to\frac4{729}.
}
\]

有限见证：
\[
n=208049,\qquad r=47,\qquad d=235.
\]
\(47\) 与六个位置
\[
208049,\ 208519,\ 209459,\ 209929,\ 210869,\ 211339
\]
全部为素数，前三组深度为 \((3,4,2)\)。

\section{文献边界}
二次多项式 \(p\)-进赋值树的两条无限分支已有一般分类；随机 \(p\)-进多项式根的分布与相关性也已有独立研究。本文不把“双根”本身视为新发现，所固化的是这个特定 \(F_3\) 平移乘积过程的显式根坐标、联合律、可逆重建和全素数局部障碍。

ROOT-1/2/3 保留既有 \texttt{PAPER-AUDITED} 状态。ROOT-4 的全族比例于 2026 年 9 月 30 日修正为 $4/729$，状态为 \texttt{PAPER-PROVED}，修订稿复核待完成；尚未 Lean 形式化。

\begin{thebibliography}{9}
\bibitem{BHS} Will Boultinghouse et al., \emph{The \(p\)-Adic Valuation Trees for Quadratic Polynomials for Odd Primes}, arXiv:2309.16637.
\bibitem{Caruso} Xavier Caruso, \emph{Where are the zeroes of a random \(p\)-adic polynomial?}, arXiv:2110.03942.
\bibitem{GT10} Ben Green and Terence Tao, \emph{Linear equations in primes}, Ann. of Math. 171 (2010), 1753--1850.
\bibitem{GTZ12} Ben Green, Terence Tao and Tamar Ziegler, \emph{An inverse theorem for the Gowers \(U^{s+1}[N]\)-norm}, Ann. of Math. 176 (2012), 1231--1372.
\end{thebibliography}
\end{document}
